% ----------------------------------
%   Calibration, numerical experiments (Chapter 7)
% ----------------------------------

\begin{frame}
\frametitle{Example 1. Constant volatility}
\begin{scriptsize}
Egger and Engl, Example 1: recover $a^\dagger = 0.15$ ($\lvol \approx 0.548$) on $y \in [-3,3]$, $\tau \in [0,1]$,
$r = q = 0$, $S_0 = 1000$, prior $a^*(y) = 0.15 - 0.05\,\operatorname{erf}(-y^2)$, $\beta = 10^{-6}$.
Cases: (A) exact data, (B) 20 strikes, (C) poor prior $a^* \equiv 0.15$, (D) $0.1\%$ noise, (BD) both.
\end{scriptsize}
\begin{center}
\begin{tiny}
\begin{tabular}{c c c c c c c}
\toprule
Strike & True & A & B & C & D & BD \\
\midrule
600  & 0.1500 & 0.1493 & 0.1500 & 0.1519 & 0.1378 & 0.1245 \\
1000 & 0.1500 & 0.1499 & 0.1499 & 0.1515 & 0.1490 & 0.1455 \\
1400 & 0.1500 & 0.1497 & 0.1498 & 0.1504 & 0.1489 & 0.1522 \\
1600 & 0.1500 & 0.1508 & 0.1501 & 0.1491 & 0.1543 & 0.1500 \\
1800 & 0.1500 & 0.1487 & 0.1504 & 0.1511 & 0.1461 & 0.1463 \\
\bottomrule
\end{tabular}
\end{tiny}
\end{center}
\begin{figure}[htb!]
    \centering
    \includegraphics[width=0.8\linewidth]{egger_example1_caseBD}
    \caption{\scriptsize Case BD: incomplete and noisy data. Good recovery where observed, prior dominates the wings.}
\end{figure}
\end{frame}

\begin{frame}
\frametitle{Example 2. The role of the prior}
\begin{scriptsize}
Oscillatory true coefficient near the money,
$a^\dagger(y) = 0.15 + 0.05\,e^{-(y+0.3)^2}\sin(2\pi y) + 0.05\,\operatorname{erf}(20 y)$,
prior $a^*(y) = 0.15 + 0.05\,\operatorname{erf}(2y)$.
Data generated on a finer mesh and wider domain to avoid the \emph{inverse crime};
noisy cases use the $\beta$ schedule with the discrepancy principle.
\end{scriptsize}
\begin{figure}[htb!]
    \centering
    \includegraphics[width=0.9\linewidth]{egger_example2_recovery}
    \caption{\scriptsize Recovered $a(y)$ (blue), true $a^\dagger$ (red dashed), prior (grey).
    In case C the constant prior cannot recover the wings, where the data are weak.}
\end{figure}
\end{frame}

\begin{frame}
\frametitle{Market surface. Extension to time-dependent volatility}
\begin{scriptsize}
\begin{itemize}
\item Prices are observed at finitely many maturities, so $a$ is taken \textbf{piecewise constant in time},
\begin{equation*}
    a(y,\tau) = \sum_{n=1}^{K} a_n(y)\,\chi_{(\tau_{n-1},\,\tau_n]}(\tau),
    \qquad a_n \in H^1(-M,M),
\end{equation*}
one slice per quoted maturity. $K = 1$ is the setting of Egger and Engl.

\item The functional adds a penalty on jumps between slices,
\begin{equation*}
    J(a) = \sum_{j=1}^{N} \bigl\| u(a)(\cdot,\tau_j) - u^\delta_j \bigr\|^2_{L^2(\Omega_j)}
         + \beta_y \sum_{n=1}^{K} \| a_n - a^*_n \|^2_{H^1(\Omega)}
         + \beta_\tau \sum_{n=2}^{K} \| a_n - a_{n-1} \|^2_{L^2(\Omega)}.
\end{equation*}
One forward march to the last maturity and one adjoint sweep give the whole gradient.

\item \textbf{Prior.} Seeding with $\tfrac12\ivol^2$ anchors the slope at half the right value (rule of two).
Instead invert the short maturity harmonic mean relation of Berestycki, Busca and Florent,
\begin{equation*}
    \frac{y}{\ivol(y)} = \int_0^y \frac{dx}{\lvol(x)}
    \quad \Longrightarrow \quad
    \lvol(y) = \frac{\ivol(y)}{1 - y\,\partial_y\ivol(y)/\ivol(y)}.
\end{equation*}
\end{itemize}
\medskip
Berestycki, H., Busca, J., Florent, I. (2002). \emph{Asymptotics and calibration of local volatility models}. Quantitative Finance.
\end{scriptsize}
\end{frame}

\begin{frame}
\frametitle{Market surface. SPY, 30 April 2026}
\begin{scriptsize}
$S_0 = 720.65$, eleven maturities $\tau \in [0.10, 0.88]$, liquid band $0.8 \le K/S_0 \le 1.2$,
$\beta_y = 0.1$, $\beta_\tau = 1$.
\end{scriptsize}
\begin{figure}[htb!]
    \centering
    \begin{subfigure}[b]{0.5\linewidth}\centering
    \includegraphics[width=\linewidth]{egger_example4_surface_3d}
    \caption{\scriptsize Calibrated $\lvol(y,\tau)$.}
    \end{subfigure}%
    \begin{subfigure}[b]{0.5\linewidth}\centering
    \includegraphics[width=\linewidth]{egger_example4_surface_slices}
    \caption{\scriptsize Local (solid) vs.\ implied (faint) volatility.}
    \end{subfigure}
\end{figure}
\begin{scriptsize}
Pronounced negative skew: $\lvol \approx 0.15$ ATM at the short end to $0.23$ at the long end,
up to $0.95$ on the short downside.
Median reprice error $0.13\%$, below $1\%$ on liquid strikes at every maturity
(the naive prior $\tfrac12\ivol^2$ gives $0.20\%$ and maxima above $2\%$).
\end{scriptsize}
\end{frame}

\begin{frame}
\frametitle{Market surface. Free of arbitrage by construction}
\begin{scriptsize}
\begin{itemize}
\item The optimizer keeps $a \geq 10^{-8}$ and the prices solve the Dupire equation with that coefficient,
so they are the prices of a genuine diffusion: the conditions of the no-arbitrage slide cannot fail.
\item The implied volatility interpolation route carries no such guarantee.
\end{itemize}
\end{scriptsize}
\begin{figure}[htb!]
    \centering
    \includegraphics[width=\linewidth]{egger_example7_arbitrage}
    \caption{\scriptsize Left: risk-neutral density $\partial^2_{KK}\mathcal{C} \geq 0$ per maturity (no butterfly arbitrage).
    Right: total implied variance $\ivol^2\tau$ increasing in maturity (no calendar arbitrage).}
\end{figure}
\end{frame}

\begin{frame}
\frametitle{Market surface. The rule of two}
\begin{scriptsize}
Near the money $\partial_y \lvol|_{y=0} = 2\,\partial_y \ivol|_{y=0}$.
Quadratic fits on $|y| \le 0.05$ of the market smile and of the calibrated slices:
\end{scriptsize}
\begin{figure}[htb!]
    \centering
    \includegraphics[width=0.8\linewidth]{egger_example5_skew}
\end{figure}
\begin{center}
\begin{tiny}
\setlength{\tabcolsep}{4pt}
\begin{tabular}{c c c c c c c c c c c c}
\toprule
$\tau$ & 0.096 & 0.132 & 0.164 & 0.211 & 0.249 & 0.307 & 0.384 & 0.416 & 0.633 & 0.710 & 0.882 \\
\midrule
ratio & 2.01 & 2.17 & 2.02 & 2.36 & 2.24 & 2.32 & 2.19 & 2.24 & 2.47 & 2.17 & 2.38 \\
\bottomrule
\end{tabular}
\end{tiny}
\end{center}
\begin{scriptsize}
The ratio averages $2.16$ below three months and stays in $[2.0, 2.5]$: the theory is confirmed empirically.
\end{scriptsize}
\end{frame}

\begin{frame}
\frametitle{Market surface. Out of sample pricing}
\begin{scriptsize}
Remove the December 2026 expiration ($\tau = 0.63$) and price it out of sample:
\begin{itemize}
\item \textbf{Local volatility}: calibrate on the other ten maturities, the Dupire equation propagates across the gap.
\item \textbf{Implied volatility}: interpolate the total variance $\ivol^2\tau$ linearly between $\tau = 0.42$ and $0.71$.
\end{itemize}
\end{scriptsize}
\begin{figure}[htb!]
    \centering
    \includegraphics[width=0.55\linewidth]{egger_example6_holdout}
\end{figure}
\begin{scriptsize}
Neither is off by more than $2\%$ on $|y| \le 0.15$.
Median error: $0.75\%$ local volatility, $0.29\%$ interpolation.
the interpolation is built to reproduce the smile, while the local volatility error accumulates through the PDE.
Yet only the local volatility surface is guaranteed to be arbitrage-free.
\end{scriptsize}
\end{frame}
